\section{Exact certificates for the Toeplitz BW form at small orders}
\label{sec:exact-certificates}

We have shown that, for sufficiently large $N$, the general Toeplitz
B\"ottcher--Wenzel form is not SoS.  In the previous section, we also
showed that several structured subclasses remain SoS at every order.

This naturally raises the question of what happens for the small orders
that arise in applications, when no additional symmetry structure is
assumed.  Is the unrestricted Toeplitz BW form still SoS?
For small orders, the answer is positive.  In this section, we show
that the unrestricted Toeplitz BW form is SoS for every
$
2\le N\le50.
$
Moreover, for
$
2\le N\le20,
$
we provide a Lean-verified proof of the SoS property\footnote{Note that the orders $21$--$50$ are certified by exact arithmetic but are
not included in the Lean verification.}.

Together with the large-order negative result, this leaves a substantial
gap.  The quantitative argument in Appendix~\ref{app:quantitative}
gives the explicit sufficient bound
$
N_0=\SufficientOrder
$
for non-SoS representability, whereas the exact positive construction
below has currently been verified only through order $50$.  The latter
stops because the positive increment used in the induction has only
been constructed and verified through the step $n=49$; this does not
imply that $F_{51}$ is not SoS.  Thus the exact SoS status of the
unrestricted Toeplitz BW form for
$
50<N<N_0
$
remains open.  Determining the first non-SoS order, or substantially
narrowing this gap, is a natural direction for future work.

\paragraph{Toeplitz BW SoS for $2\le N\le50$}\label{subsec:certificate-chain}
We establish the result by an exact induction. Starting from $F_2$,
each step expresses the next Toeplitz BW form as a positive multiple
of a linear substitution of the preceding form, plus an SoS increment.

Here $N$ denotes the matrix order in the final result, while $n$ is
used only as the induction index. At the step from $n$ to $n+1$, let
$z^{(n+1)}$ denote the wedge vector associated with the symbol variables
of $F_{n+1}$.

\begin{lemma}[Exact positive increments]
\label{lem:finite-increments}
For each $n=2,\ldots,49$, the supplied rational data specify a
positive rational number $\alpha_n$, a rational linear map
$\psi_n:\R^{4n+2}\to\R^{4n-2}$, and a rational symmetric matrix
$\mathbf H_n\succeq0$ such that
\begin{equation}\label{eq:finite-chain}
 F_{n+1}
 =
 \alpha_n(F_n\circ\psi_n)
 +(z^{(n+1)})^\top\mathbf H_nz^{(n+1)}.
\end{equation}
\end{lemma}

\begin{proof}
Each identity is verified by reconstructing the two quartics from
their defining Toeplitz matrices and comparing all polynomial
coefficients over $\mathbb Q$. Positive semidefiniteness of
$\mathbf H_n$ is checked by exact rational matrix calculations after
removing explicitly verified kernel directions. The data for all
$48$ steps and the verification procedure are given in
Appendix~\ref{app:finite-certificates}.
\end{proof}


\begin{theorem}[The finite positive range]
\label{thm:finite-positive}
For every integer $2\le N\le50$, $F_N$ is a sum of squares of real
homogeneous quadratic polynomials. For $2\le N\le20$, the same
statement is also formally verified in Lean~4.
\end{theorem}

\begin{proof}
At order two, direct expansion gives
$
 F_2=4z_{-1,0}^2+4z_{0,1}^2.
$
Suppose $F_n=\sum_j q_j^2$, where the $q_j$ are real homogeneous
quadratics. Since $\mathbf H_n\succeq0$, it admits a real factorization
$\mathbf H_n=V_n^\top V_n$. Equation~\eqref{eq:finite-chain} gives
\[
 F_{n+1}
 =
 \sum_j\bigl(\sqrt{\alpha_n}\,q_j\circ\psi_n\bigr)^2
 +
 \sum_i\bigl(V_nz^{(n+1)}\bigr)_i^2.
\]
Each factor is again a homogeneous quadratic. Induction through
$n=49$ therefore proves the result for every $2\le N\le50$.
\end{proof}

The substitutions in \eqref{eq:finite-chain} allow the SoS form at one
order to be reused at the next. For $2\le n\le16$, a simple rescaling
of $F_n$ is sufficient. For $17\le n\le49$, the substitution also
damps the symbol at offset $a$ by a rational factor $q_n^{|a|}$,
with $0<q_n\le1$, so that the remaining increment admits a positive
semidefinite Gram representation. The precise constructions are given
in Appendix~\ref{app:finite-certificates}. 

\paragraph{Lean verification through order $20$.} For $2\le N\le20$, the SoS assertion is also formally verified in
Lean~4.  In particular, the range $13\le N\le20$ extends the previously
known positive range through order $12$.  The Lean formalization works
with the original Toeplitz BW polynomial and proves directly that it is
a finite sum of real homogeneous quadratic squares.  The corresponding
formal statement and source dependencies are recorded in
Appendix~\ref{app:positive-lean}.

This verification is distinct from the exact rational induction in
\Cref{subsec:certificate-chain}.  The rational argument proves the SoS
property for every $2\le N\le50$ by checking the identities
\eqref{eq:finite-chain} and the positive semidefiniteness of the
associated rational matrices.  The Lean proof, on the other hand,
provides an independent formal verification for each order
$2\le N\le20$ and need not use the same Gram representations. 
Thus,
\[
 \begin{array}{c|c}
 \text{orders}&\text{verification}\\\hline
 2\le N\le20&\text{exact rational verification and Lean~4}\\
 21\le N\le50&\text{exact rational verification}
 \end{array}
\]
Both results concern unrestricted real Toeplitz pairs.  The status of
the unrestricted form remains open for the intermediate range
$
50<N<N_0.
$